%% ----------------------------------------------------------------------------
%%  §7.3  Derivate parziali. Gradiente
%% ----------------------------------------------------------------------------
\section{Derivate parziali. Gradiente}
\label{sec:derivate-parziali}

Sia $f$ una funzione di due variabili definita in un intorno circolare $I_{\delta}$ di un
punto $(x,y)$. In particolare $f$ è definita, oltre che in $(x,y)$, anche nei punti
$(x+h,y)$ con $\abs h<\delta$, che stanno sulla retta orizzontale per $(x,y)$, e nei punti
$(x,y+k)$ con $\abs k<\delta$, sulla retta verticale. Fermando una delle due variabili si
ottiene una funzione di una sola variabile, che possiamo derivare.

\begin{definizione}[Derivate parziali]\label{def:derivate-parziali}\index{derivata!parziale}
La \emph{derivata parziale di $f$ rispetto a $x$} nel punto $(x,y)$ è il limite
\begin{equation}\label{eq:derivata-parziale-x}
  \lim_{h\to0}\frac{f(x+h,y)-f(x,y)}{h},
\end{equation}
se esiste ed è finito; si indica con uno dei simboli
\begin{equation}\label{eq:notazioni-parziale-x}
  f_{x}, \qquad f_{x}(x,y), \qquad \frac{\partial f}{\partial x}, \qquad D_{x}f .
\end{equation}
Analogamente, la \emph{derivata parziale di $f$ rispetto a $y$} nel punto $(x,y)$ è il
limite
\begin{equation}\label{eq:derivata-parziale-y}
  \lim_{k\to0}\frac{f(x,y+k)-f(x,y)}{k},
\end{equation}
se esiste ed è finito, e si indica con $f_{y}$, $f_{y}(x,y)$, $\frac{\partial f}{\partial y}$,
$D_{y}f$.
\end{definizione}

Si noti che la \eqref{eq:derivata-parziale-x} è il limite del rapporto incrementale di una
funzione di \emph{una} variabile: posto $g(t)=f(t,y)$, con $y$ fissato che gioca il ruolo
di parametro, è $f_{x}(x,y)=g'(x)$. Per questo motivo \emph{le derivate parziali si
calcolano con le stesse regole di derivazione delle funzioni di una variabile,
considerando l'altra variabile costante}. Geometricamente, $f_{x}(x_{0},y_{0})$ è la
pendenza della curva che si ottiene sezionando il grafico con il piano verticale
$y=y_{0}$, e $f_{y}(x_{0},y_{0})$ quella della sezione con il piano $x=x_{0}$
(\cref{fig:derivate-parziali}).

\begin{figure}[htb]
\centering
\includegraphics[width=0.62\textwidth]{derivate_parziali}
\caption{Il significato geometrico delle derivate parziali per
$f(x,y)=3-\frac{x^{2}}2-\frac{y^{2}}4$ nel punto $(x_{0},y_{0})=(1,1)$. La sezione con il
piano $y=y_{0}$ (rame) ha nel punto $(x_{0},y_{0},f(x_{0},y_{0}))$ tangente di pendenza
$f_{x}(1,1)=-1$; la sezione con il piano $x=x_{0}$ (verde) ha tangente di pendenza
$f_{y}(1,1)=-\frac12$ (\texttt{python/fig\_derivate\_parziali.py}).}
\label{fig:derivate-parziali}
\end{figure}

\begin{esempio}[Calcolo di derivate parziali]\label{es:derivate-parziali}
Per la funzione
\begin{equation}\label{eq:parziali-es1}
  f(x,y)=x^{2}+xy^{2}, \qquad (x,y)\in\R^{2},
\end{equation}
nel derivare rispetto a $x$ si considera $y$ costante, come se fosse, ad esempio,
$y^{2}=9$ e si dovesse derivare la funzione $x\mapsto x^{2}+9x$; si ottiene
$f_{x}(x,y)=2x+y^{2}$. Analogamente, derivando rispetto a $y$ con $x$ costante,
$f_{y}(x,y)=2xy$. La funzione $f(x,y)=\sen(xy)$, il cui grafico è nella
\cref{fig:superfici}, ha derivate parziali $f_{x}=y\cos(xy)$ e $f_{y}=x\cos(xy)$ in ogni
punto di $\R^{2}$ (regola della catena). La funzione $f(x,y)=e^{x}$ è costante rispetto a
$y$: le sue derivate parziali sono $f_{x}=e^{x}$ e $f_{y}=0$.
\end{esempio}

\begin{python}[Derivate parziali con sympy e controllo numerico]
import sympy as sp

x, y = sp.symbols("x y", real=True)
for f in (x**2 + x*y**2, sp.sin(x*y), sp.exp(x), x + 2*y):
    gx, gy = sp.diff(f, x), sp.diff(f, y)
    print(f"f = {str(f):13}  f_x = {str(gx):10}  f_y = {gy}")

# controllo numerico con i rapporti incrementali, nel punto (1, 2)
F = sp.lambdify((x, y), x**2 + x*y**2)
h = 1e-6
print("f_x(1,2) ~", round((F(1 + h, 2) - F(1, 2)) / h, 4),
      "  f_y(1,2) ~", round((F(1, 2 + h) - F(1, 2)) / h, 4))
\end{python}
\begin{risultato}
f = x**2 + x*y**2  f_x = 2*x + y**2  f_y = 2*x*y
f = sin(x*y)       f_x = y*cos(x*y)  f_y = x*cos(x*y)
f = exp(x)         f_x = exp(x)      f_y = 0
f = x + 2*y        f_x = 1           f_y = 2
f_x(1,2) ~ 6.0   f_y(1,2) ~ 4.0
\end{risultato}

\begin{esempio}[Sensibilità della resistenza all'allungamento alare]\label{es:sensibilita-AR}
Con la polare parabolica, il coefficiente di resistenza dipende dal coefficiente di
portanza e dall'allungamento alare:
\begin{equation}\label{eq:CD-CL-AR}
  C_{D}(C_{L},AR)=C_{D0}+\frac{C_{L}^{2}}{\pi e\,AR} .
\end{equation}
Le derivate parziali misurano la \emph{sensibilità} di $C_{D}$ a ciascuna delle due
variabili, a parità dell'altra:
\[
  \frac{\partial C_{D}}{\partial C_{L}}=\frac{2C_{L}}{\pi e\,AR}, \qquad
  \frac{\partial C_{D}}{\partial AR}=-\frac{C_{L}^{2}}{\pi e\,AR^{2}} .
\]
La seconda è negativa: aumentando l'allungamento la resistenza indotta diminuisce, ma
sempre meno, perché il modulo della derivata decresce come $\frac1{AR^{2}}$; ed è tanto
maggiore in modulo quanto più alto è $C_{L}$, cioè nel volo lento, in salita e in
virata. È il motivo per cui alianti e velivoli da alta quota hanno ali molto allungate.
\end{esempio}

\begin{python}[Sensibilità di $C_{D}$ a $C_{L}$ e ad $AR$]
import math

# C_D(C_L, AR) = C_D0 + C_L^2 / (pi e AR): sensibilità parziali
CD0, e = 0.025, 0.8
CD = lambda CL, AR: CD0 + CL**2 / (math.pi * e * AR)
dCD_dCL = lambda CL, AR: 2 * CL / (math.pi * e * AR)
dCD_dAR = lambda CL, AR: -CL**2 / (math.pi * e * AR**2)
print("  C_L  AR     C_D     dC_D/dC_L   dC_D/dAR")
for CL, AR in ((0.5, 8), (1.0, 8), (1.0, 12)):
    print(f"  {CL:.1f}  {AR:2d}   {CD(CL, AR):.5f}   "
          f"{dCD_dCL(CL, AR):.5f}    {dCD_dAR(CL, AR):.6f}")
\end{python}
\begin{risultato}
  C_L  AR     C_D     dC_D/dC_L   dC_D/dAR
  0.5   8   0.03743   0.04974    -0.001554
  1.0   8   0.07474   0.09947    -0.006217
  1.0  12   0.05816   0.06631    -0.002763
\end{risultato}

\subsection{Il gradiente e le derivate direzionali}

\begin{definizione}[Gradiente]\label{def:gradiente}\index{gradiente}
Se $f$ ammette le derivate parziali $f_{x}$, $f_{y}$ nel punto $(x,y)$, si chiama
\emph{gradiente} di $f$ in $(x,y)$ il vettore di $\R^{2}$ che ha per componenti le
derivate parziali:
\begin{equation}\label{eq:gradiente}
  \nabla f(x,y)=\operatorname{grad}f(x,y)=\bigl(f_{x}(x,y),\,f_{y}(x,y)\bigr).
\end{equation}
\end{definizione}

Il simbolo $\nabla$ si legge «nabla» (dal nome di un'arpa fenicia, di forma triangolare).
Le derivate parziali descrivono come varia $f$ quando ci si sposta parallelamente agli
assi; ci si può però spostare in qualunque direzione. Sia $v=(\cos\vartheta,\sen\vartheta)$
un \emph{versore}, cioè un vettore di lunghezza $1$; la \emph{derivata direzionale} di $f$
nel punto $P=(x,y)$ secondo la direzione $v$ è il limite
\begin{equation}\label{eq:derivata-direzionale}
  \frac{\partial f}{\partial v}(x,y)=\lim_{t\to0}\frac{f(x+t\cos\vartheta,\,y+t\sen\vartheta)-f(x,y)}{t},
\end{equation}
se esiste finito. Per $\vartheta=0$ si ritrova $f_{x}$, per $\vartheta=\frac\pi2$ si ritrova
$f_{y}$. Il gradiente ha una proprietà notevole: \emph{se non è nullo, indica la direzione
di massima pendenza del grafico}. Lo verifichiamo qui per le funzioni lineari; il caso
generale, per le funzioni differenziabili, è trattato nel \cref{sec:differenziabilita}.

\begin{proposizione}[Gradiente di una funzione lineare]\label{prop:gradiente-lineare}
Sia $f(x,y)=ax+by+c$, con $(a,b)\neq(0,0)$. Allora $\nabla f=(a,b)$ in ogni punto, e per
ogni versore $v=(\cos\vartheta,\sen\vartheta)$
\begin{equation}\label{eq:direzionale-lineare}
  \frac{\partial f}{\partial v}=a\cos\vartheta+b\sen\vartheta=\abs{\nabla f}\cos(\vartheta-\varphi),
\end{equation}
dove $\abs{\nabla f}=\sqrt{a^{2}+b^{2}}$ e $\varphi$ è l'angolo che il gradiente forma con
l'asse $x$. In particolare la derivata direzionale è massima, e uguale a $\abs{\nabla f}$,
nella direzione del gradiente; è minima nella direzione opposta; è nulla nelle direzioni
perpendicolari al gradiente, che sono quelle delle curve di livello.
\end{proposizione}

\begin{schema}
\Ipotesi{$f(x,y)=ax+by+c$, $(a,b)\neq(0,0)$; $v=(\cos\vartheta,\sen\vartheta)$.}
\Tesi{$\frac{\partial f}{\partial v}=\abs{\nabla f}\cos(\vartheta-\varphi)$: massima per
$v$ parallelo e concorde a $\nabla f$, nulla per $v$ perpendicolare.}
\end{schema}

\begin{strategia}
Per una funzione lineare il rapporto incrementale è costante, e la derivata direzionale
si calcola subito. Scrivendo il gradiente in forma polare,
$(a,b)=\abs{\nabla f}(\cos\varphi,\sen\varphi)$, la formula di sottrazione del coseno
trasforma $a\cos\vartheta+b\sen\vartheta$ in un solo coseno, di cui si conoscono massimo e
minimo.
\end{strategia}

\begin{dimostrazione}
Si ha $f_{x}=a$ e $f_{y}=b$. Inoltre
$f(x+t\cos\vartheta,y+t\sen\vartheta)-f(x,y)=t(a\cos\vartheta+b\sen\vartheta)$, quindi il rapporto
incrementale in \eqref{eq:derivata-direzionale} è costante e uguale a
$a\cos\vartheta+b\sen\vartheta$. Poiché $(a,b)\neq(0,0)$, posto $r=\sqrt{a^{2}+b^{2}}$ esiste
$\varphi$ con $a=r\cos\varphi$ e $b=r\sen\varphi$; per la formula di sottrazione
\eqref{eq:addizione-cos},
\[
  a\cos\vartheta+b\sen\vartheta=r(\cos\varphi\cos\vartheta+\sen\varphi\sen\vartheta)=r\cos(\vartheta-\varphi).
\]
Il massimo $r$ si ha per $\vartheta=\varphi$, il minimo $-r$ per $\vartheta=\varphi+\pi$, e il
valore $0$ per $\vartheta=\varphi\pm\frac\pi2$. Infine, le curve di livello
$ax+by+c=\text{costante}$ sono rette di direzione $(-b,a)$, perpendicolare a $(a,b)$.
\end{dimostrazione}

Per esempio, la funzione $f(x,y)=x+2y$ ha gradiente costante $(1,2)$: questo vettore
indica la direzione e il verso, nel piano di base, in cui conviene muoversi per ottenere
il massimo incremento di $f$ a parità di percorso. Il grafico di una funzione lineare è
un piano; per una funzione qualunque, nel \cref{sec:differenziabilita} vedremo che
vicino a ogni punto il grafico si confonde con un piano (il \emph{piano tangente}), e
che per questo la proprietà vale in generale.

\begin{idea}[La forza di gradiente di pressione]
Su una particella d'aria agisce, per effetto delle differenze di pressione, una forza
per unità di massa uguale a $-\frac1\rho\nabla p$: la forza è diretta come l'opposto del
gradiente, cioè dalle alte verso le basse pressioni, perpendicolarmente alle isobare, ed è
tanto più intensa quanto più le isobare sono fitte (\cref{fig:livello}, a sinistra). Se la
Terra non ruotasse, il vento soffierebbe proprio così; nella realtà la forza di Coriolis
lo devia, e alle quote alte il vento scorre quasi parallelo alle isobare (vento
geostrofico). La stessa forza governa lo strato limite sul profilo di un'ala: dove la
pressione cresce nel verso del moto (\emph{gradiente avverso}) la corrente rallenta e può
separarsi dalla superficie, provocando lo stallo.
\end{idea}

\begin{notastorica}[La nascita delle derivate parziali e del simbolo $\partial$]
Le derivate parziali nacquero insieme al calcolo, ma si affermarono negli anni Trenta e
Quaranta del Settecento, nei lavori di Eulero, di Clairaut e di d'Alembert sulla
meccanica dei corpi continui: nel 1747 d'Alembert scrisse l'equazione della corda
vibrante, la prima equazione alle derivate parziali della fisica matematica. Nessuno,
però, si fermò a darne una definizione esplicita: le si calcolava trattando come
costanti le altre variabili, esattamente come facciamo noi, e soltanto la revisione
critica dell'Ottocento rivelò che cosa esse non dicono
(\cref{sec:differenziabilita}). Per molto tempo si usò lo stesso simbolo $d$ delle
derivate ordinarie, specificando a parole la variabile. Il simbolo $\partial$ compare
già in una memoria di Condorcet del 1770; nella forma $\frac{\partial u}{\partial x}$ fu
usato da Adrien-Marie Legendre nel 1786, poi abbandonato, e reintrodotto in modo
definitivo da Carl Gustav Jacobi nel 1841. Il simbolo $\nabla$ e il nome «nabla» si
affermarono nell'Ottocento con il calcolo vettoriale di Hamilton, Tait, Maxwell e
Gibbs (il nome fu suggerito a Tait, intorno al 1870, dal biblista William Robertson
Smith).
\end{notastorica}
